\documentclass{article}
\usepackage{amsmath,amssymb,amsthm}
\usepackage{algorithm,algorithmic}
\usepackage{enumitem}
\usepackage{hyperref}
\usepackage{geometry}
\geometry{margin=1in}
\newcommand{\E}{\mathbb{E}}
\newcommand{\R}{\mathbb{R}}
\newcommand{\norm}[1]{\left\lVert#1\right\rVert}
\newtheorem{theorem}{Theorem}
\newtheorem{lemma}{Lemma}
\newtheorem{assumption}{Assumption}

\begin{document}

\section*{Probabilistic Coordinate Descent (PCD)}

\section*{Mathematical Formulation}
Let $f:\R^{d}\to\R$ be differentiable.  
At iteration $k$ a random coordinate $i_{k}\in\{1,\dots ,d\}$ is drawn
according to a fixed probability vector $p=(p_{1},\dots ,p_{d})$,
\[
\Pr(i_{k}=i)=p_{i},\qquad p_{i}>0,\;\sum_{i=1}^{d}p_{i}=1 .
\]

Denote by $e_{i}$ the $i$‑th canonical basis vector.
Given a stepsize $\eta_{k}>0$, the update of the $i_{k}$‑th coordinate is
\[
\boxed{%
w_{k+1}=w_{k}-\frac{\eta_{k}}{p_{i_{k}}}\,e_{i_{k}}\,\partial_{i_{k}}f(w_{k})},
\tag{1}
\]
where $\partial_{i}f(w)=\frac{\partial f(w)}{\partial w^{(i)}}$.
All other coordinates remain unchanged.

\section*{Pseudocode}
\begin{algorithm}[H]
\caption{Probabilistic Coordinate Descent}
\begin{algorithmic}[1]
\STATE {\bf Input:} initial point $w_{0}\in\R^{d}$, stepsizes $\{\eta_{k}\}_{k\ge0}$, probability vector $p$.
\FOR{$k=0,1,2,\dots$}
    \STATE Sample $i_{k}\sim p$.
    \STATE Compute partial gradient $g_{k}\leftarrow \partial_{i_{k}}f(w_{k})$.
    \STATE Update 
    \[
    w_{k+1}\leftarrow w_{k}-\frac{\eta_{k}}{p_{i_{k}}}\,e_{i_{k}}\,g_{k}.
    \]
    \IF{termination criterion met (e.g. $\|w_{k+1}-w_{k}\|\le\epsilon$)}
        \STATE \textbf{break}
    \ENDIF
\ENDFOR
\STATE {\bf Output:} $w_{K}$ (last iterate) or $\tilde w=\frac1K\sum_{k=0}^{K-1}w_{k}$.
\end{algorithmic}
\end{algorithm}

\section*{Dry Run Example}
Consider the scalar function $f(w)=(w-3)^{2}$, thus $d=1$,
$\partial_{1}f(w)=2(w-3)$.
Choose $p_{1}=1$, stepsize $\eta=0.1$, and $w_{0}=0$.

\begin{enumerate}[label=Iteration \arabic*:]
\item \textbf{Gradient:} $g_{0}=2(w_{0}-3)=2(-3)=-6$.
    
    \textbf{Update:}
    \[
    w_{1}=w_{0}-\frac{0.1}{1}\,g_{0}=0-0.1(-6)=0.6 .
    \]
    \textbf{Objective:} $f(w_{1})=(0.6-3)^{2}=5.76$.

\item \textbf{Gradient:} $g_{1}=2(w_{1}-3)=2(-2.4)=-4.8$.
    
    \textbf{Update:}
    \[
    w_{2}=w_{1}-0.1\,g_{1}=0.6-0.1(-4.8)=1.08 .
    \]
    \textbf{Objective:} $f(w_{2})=(1.08-3)^{2}=3.6864$.

\item \textbf{Gradient:} $g_{2}=2(w_{2}-3)=2(-1.92)=-3.84$.
    
    \textbf{Update:}
    \[
    w_{3}=w_{2}-0.1\,g_{2}=1.08-0.1(-3.84)=1.464 .
    \]
    \textbf{Objective:} $f(w_{3})=(1.464-3)^{2}=2.366 .
    \]
\end{enumerate}
The objective value decreases at each step, illustrating the algorithm.

\newpage
\section*{Assumptions}
\begin{assumption}[Coordinate‑wise $L_{i}$‑smoothness]\label{ass:smooth}
For each coordinate $i\in\{1,\dots ,d\}$ there exists $L_{i}>0$ such that for all $w\in\R^{d}$ and any $\alpha\in\R$,
\[
f(w+\alpha e_{i})\le f(w)+\alpha\,\partial_{i}f(w)+\frac{L_{i}}{2}\alpha^{2}.
\tag{2}
\]
\end{assumption}

\begin{assumption}[Bounded below]\label{ass:lower}
$f$ is bounded from below: $f^{\inf}\triangleq\inf_{w}f(w)>-\infty$.
\end{assumption}

\begin{assumption}[Stepsize]\label{ass:stepsize}
The stepsizes satisfy $0<\eta_{k}\le\displaystyle\frac{1}{2L_{\max}}$,
where $L_{\max}\triangleq\max_{i}L_{i}$.
\end{assumption}

\section*{Main Convergence Theorem}
\begin{theorem}[Expected gradient norm decay]\label{thm:main}
Let $\{w_{k}\}$ be generated by \eqref{1} under Assumptions~\ref{ass:smooth}–\ref{ass:stepsize}.
Define the weighted norm $\|g\|_{p^{-1}}^{2}\triangleq\sum_{i=1}^{d}\frac{g_{i}^{2}}{p_{i}}$.
Then for any $K\ge1$,
\[
\frac{1}{K}\sum_{k=0}^{K-1}\E\!\left[\|\nabla f(w_{k})\|^{2}\right]
\;\le\;
\frac{2L_{\max}}{\eta_{\min}}\,
\frac{f(w_{0})-f^{\inf}}{K}\,
\max_{i}p_{i},
\tag{3}
\]
where $\eta_{\min}\triangleq\min_{k}\eta_{k}>0$.
Consequently,
\[
\min_{0\le k\le K-1}\E\!\left[\|\nabla f(w_{k})\|^{2}\right]
=O\!\left(\frac{1}{K}\right).
\]
\end{theorem}

\section*{Theoretical Properties}
\begin{itemize}
\item \textbf{Stability.} Because the update is scaled by $1/p_{i_{k}}$, coordinates that are sampled rarely receive larger steps, preventing systematic neglect.
\item \textbf{Hyper‑parameter sensitivity.}
    \begin{enumerate}
    \item The bound \eqref{3} shows linear dependence on $\max_{i}p_{i}$; a more uniform $p$ (i.e., $p_{i}=1/d$) yields the smallest constant.
    \item The stepsize must respect $\eta_{k}\le 1/(2L_{\max})$; using a larger $\eta_{k}$ may break inequality \eqref{Step 6} below.
    \end{enumerate}
\item \textbf{Comparison to uniform random coordinate descent.}
    If $p_{i}=1/d$ for all $i$, \eqref{3} reduces to the classical $O(d/K)$ bound.
    The weighted formulation allows importance sampling (e.g., $p_{i}\propto L_{i}$) to improve the constant.
\end{itemize}

\section*{Discussion}
The theorem guarantees that even for non‑convex objectives the algorithm finds an
\emph{expected} stationary point at a rate $O(1/K)$.  The proof follows the
standard descent‑lemma argument but crucially incorporates the factor $1/p_{i}$
in the update, which appears in the expectation calculations.  The result
matches the optimal rate for stochastic first‑order methods under the same
smoothness assumptions.

\newpage
\section*{Preliminaries (Definitions and Lemmas)}
\begin{lemma}[Descent lemma for a sampled coordinate]\label{lem:descent}
Let $i$ be any coordinate and $0<\eta\le 1/(2L_{i})$.  Under
Assumption~\ref{ass:smooth},
\[
f\!\left(w-\frac{\eta}{p_{i}}e_{i}\partial_{i}f(w)\right)
\le f(w)-\frac{\eta}{2p_{i}}\bigl(\partial_{i}f(w)\bigr)^{2}.
\tag{4}
\]
\end{lemma}
\begin{proof}
Apply \eqref{2} with $\alpha=-\frac{\eta}{p_{i}}\partial_{i}f(w)$:
\[
\begin{aligned}
f\!\left(w-\frac{\eta}{p_{i}}e_{i}\partial_{i}f(w)\right)
&\le f(w)-\frac{\eta}{p_{i}}\bigl(\partial_{i}f(w)\bigr)^{2}
   +\frac{L_{i}}{2}\frac{\eta^{2}}{p_{i}^{2}}\bigl(\partial_{i}f(w)\bigr)^{2}\\[4pt]
&= f(w)-\frac{\eta}{p_{i}}\bigl(\partial_{i}f(w)\bigr)^{2}
   \left(1-\frac{L_{i}\eta}{2p_{i}}\right).
\end{aligned}
\]
Since $\eta\le 1/(2L_{i})$ and $p_{i}\le1$, we have $L_{i}\eta/(2p_{i})\le 1/2$,
hence $1-\frac{L_{i}\eta}{2p_{i}}\ge\frac12$, which yields \eqref{4}.
\end{proof}

\section*{Main Proof (Step‑by‑Step Derivation)}
We prove Theorem~\ref{thm:main}.  Throughout we denote by $I_{k}$ the random
coordinate sampled at iteration $k$.

\bigskip
\noindent\textbf{[Step 1]}  Write the one‑step progress using Lemma~\ref{lem:descent} with $\eta=\eta_{k}$ and $i=I_{k}$:
\[
f(w_{k+1})\le f(w_{k})-
\frac{\eta_{k}}{2p_{I_{k}}}\bigl(\partial_{I_{k}}f(w_{k})\bigr)^{2}.
\tag{5}
\]

\noindent\textbf{[Step 2]}  Take total expectation conditioned on $w_{k}$.
Because $\Pr(I_{k}=i\mid w_{k})=p_{i}$,
\[
\begin{aligned}
\E\!\left[f(w_{k+1})\mid w_{k}\right]
&\le f(w_{k})-
\frac{\eta_{k}}{2}\sum_{i=1}^{d}p_{i}\frac{(\partial_{i}f(w_{k}))^{2}}{p_{i}}\\[4pt]
&= f(w_{k})-
\frac{\eta_{k}}{2}\sum_{i=1}^{d}\frac{(\partial_{i}f(w_{k}))^{2}}{p_{i}}\\[4pt]
&= f(w_{k})-\frac{\eta_{k}}{2}\,\|\nabla f(w_{k})\|_{p^{-1}}^{2}.
\end{aligned}
\tag{6}
\]

\noindent\textbf{[Step 3]}  Remove the conditioning by taking full expectation:
\[
\E\!\left[f(w_{k+1})\right]\le \E\!\left[f(w_{k})\right]
-\frac{\eta_{k}}{2}\,\E\!\left[\|\nabla f(w_{k})\|_{p^{-1}}^{2}\right].
\tag{7}
\]

\noindent\textbf{[Step 4]}  Rearrange \eqref{7} to isolate the weighted gradient term:
\[
\frac{\eta_{k}}{2}\,\E\!\left[\|\nabla f(w_{k})\|_{p^{-1}}^{2}\right]
\le \E\!\left[f(w_{k})-f(w_{k+1})\right].
\tag{8}
\]

\noindent\textbf{[Step 5]}  Sum \eqref{8} over $k=0,\dots ,K-1$:
\[
\frac{1}{2}\sum_{k=0}^{K-1}\eta_{k}\,
\E\!\left[\|\nabla f(w_{k})\|_{p^{-1}}^{2}\right]
\le \E\!\left[f(w_{0})-f(w_{K})\right]
\le f(w_{0})-f^{\inf}.
\tag{9}
\]

\noindent\textbf{[Step 6]}  Use $\eta_{k}\ge\eta_{\min}>0$ (Assumption~\ref{ass:stepsize}) to lower‑bound the left‑hand side:
\[
\frac{\eta_{\min}}{2}\sum_{k=0}^{K-1}
\E\!\left[\|\nabla f(w_{k})\|_{p^{-1}}^{2}\right]
\le f(w_{0})-f^{\inf}.
\tag{10}
\]

\noindent\textbf{[Step 7]}  Divide by $K$ and solve for the average weighted gradient norm:
\[
\frac{1}{K}\sum_{k=0}^{K-1}
\E\!\left[\|\nabla f(w_{k})\|_{p^{-1}}^{2}\right]
\le \frac{2\bigl(f(w_{0})-f^{\inf}\bigr)}{\eta_{\min}K}.
\tag{11}
\]

\noindent\textbf{[Step 8]}  Relate the weighted norm to the standard Euclidean norm.
Since $p_{i}>0$,
\[
\|\nabla f(w)\|^{2}
=\sum_{i=1}^{d}(\partial_{i}f(w))^{2}
\le \Bigl(\max_{i}p_{i}\Bigr)\,
\sum_{i=1}^{d}\frac{(\partial_{i}f(w))^{2}}{p_{i}}
= \Bigl(\max_{i}p_{i}\Bigr)\,\|\nabla f(w)\|_{p^{-1}}^{2}.
\tag{12}
\]

\noindent\textbf{[Step 9]}  Combine \eqref{11} and \eqref{12}:
\[
\frac{1}{K}\sum_{k=0}^{K-1}
\E\!\left[\|\nabla f(w_{k})\|^{2}\right]
\le
\max_{i}p_{i}\;
\frac{2\bigl(f(w_{0})-f^{\inf}\bigr)}{\eta_{\min}K}.
\tag{13}
\]

\noindent\textbf{[Step 10]}  Since $L_{i}\le L_{\max}$ for all $i$, the admissible stepsize
condition $\eta_{\min}\le 1/(2L_{\max})$ implies $\eta_{\min}=c/(2L_{\max})$ for some
$c\in(0,1]$.  Substituting $\eta_{\min}=c/(2L_{\max})$ into \eqref{13} yields
\[
\frac{1}{K}\sum_{k=0}^{K-1}
\E\!\left[\|\nabla f(w_{k})\|^{2}\right]
\le
\frac{2L_{\max}}{c}\,
\frac{f(w_{0})-f^{\inf}}{K}\,
\max_{i}p_{i}.
\tag{14}
\]

\noindent\textbf{[Step 11]}  Setting $c=1$ (the largest stepsize allowed) gives the
clean bound \eqref{3} stated in Theorem~\ref{thm:main}.  The final inequality
\[
\min_{0\le k\le K-1}\E\!\left[\|\nabla f(w_{k})\|^{2}\right]
\le\frac{1}{K}\sum_{k=0}^{K-1}
\E\!\left[\|\nabla f(w_{k})\|^{2}\right]
=O\!\left(\frac{1}{K}\right)
\]
follows immediately.

\hfill $\square$

\section*{Rate Derivation (With Constants)}
From \eqref{14} we obtain an explicit constant:
\[
C\;=\;\frac{2L_{\max}}{c}\,\bigl(f(w_{0})-f^{\inf}\bigr)\,\max_{i}p_{i},
\qquad
\frac{1}{K}\sum_{k=0}^{K-1}\E\!\left[\|\nabla f(w_{k})\|^{2}\right]
\le \frac{C}{K}.
\]
If the probabilities are chosen as
$p_{i}=L_{i}/\sum_{j}L_{j}$ (importance sampling), then
$\max_{i}p_{i}\le \frac{L_{\max}}{\sum_{j}L_{j}}$,
which can dramatically shrink $C$ when the $L_{i}$’s are heterogeneous.

\section*{Special Cases}
\begin{itemize}
\item \textbf{Uniform sampling.} $p_{i}=1/d$ for all $i$.
    Then $\max_{i}p_{i}=1/d$ and \eqref{3} becomes
    $\displaystyle\frac{2L_{\max}d}{\eta_{\min}K}(f(w_{0})-f^{\inf})$,
    recovering the classical $O(d/K)$ bound.
\item \textbf{Exact coordinate Lipschitz sampling.}
    Choosing $p_{i}\propto L_{i}$ yields $\max_{i}p_{i}=L_{\max}/\sum_{j}L_{j}$,
    and the bound improves to $O\!\bigl((\sum_{j}L_{j})/K\bigr)$.
\item \textbf{Deterministic cyclic coordinate descent.}
    Setting $p_{i}=1$ for a single coordinate and $0$ for others reduces \eqref{1}
    to standard (full) gradient descent on that coordinate, and the proof
    collapses to the classical smooth‑nonconvex descent lemma.
\end{itemize}

\end{document}